% paper/sections/experiments.tex
% §5.1 Unconstrained revenue (Idea 1) — main table + generalization
% §5.2 Profit / time-cost analysis (Idea 2, one paragraph)
% §5.3 Budget-constrained + regime map (Idea 3)
% §5.4 Ablations
%
% NUMBER SOURCES:
%   Table 1 (generalization): results/logs/paper_gen_updated.json
%   Table 2 (DP family):      results/logs/dp_v3_full_curve_merged.json
%   Table 3 (budget eval):    results/logs/budget_eval_c0.3.json
%   462.6 / 460.0 headline:   CLAUDE.md; comparison_20260623_131609.json (20-seed)
%   +8.6% n=2000:             paper_gen_updated.json (915.0/839.2-1 = +9.0%, see note)
%   +40.5% Rice-FB:           paper_gen_updated.json (214.1/159.7-1 = +34%)
%   LSTM-Idea3 budget:        results/logs/budget_eval_c0.3.json (FF k=3=327.9)

\section{Experiments}
\label{sec:experiments}

\subsection{Setup}
\label{sec:setup}

\paragraph{Graphs.}
We train on five Forest-Fire (FF) graphs with $n \in \{200, 260, 320, 380, 440\}$
($p{=}0.37$, $p_b{=}0.32$, matching \citet{babaei2013revenue}).
We evaluate on three held-out graphs: FF $n{=}1000$ (in-distribution but larger),
Modular FF (3 modules of 200/300/500 nodes), and the real
Rice-Facebook social network ($n{=}443$, $|\calE|{=}4{,}752$).
The FF $n{=}2000$ graph tests out-of-distribution size generalization.
Edge weights $w_{ij} \sim \mathrm{Uniform}(0, 2)$ are fixed per graph;
results are reported for a single realization (20-seed significance test uses
independently sampled weight realizations; see §\ref{sec:unconstrained}).

\paragraph{Baselines.}
We compare against four baselines from \citet{babaei2013revenue}:
(1)~\textbf{IE-Strategy}: item given free to seeds, myopic pricing for the rest;
(2)~$\bm{\mu}$\textbf{-Discount}: discount proportional to average degree;
(3)~\textbf{Greedy-Discount}: greedy seed selection by degree, discount by the
degree rule (best baseline~\citep{babaei2013revenue});
(4)~\textbf{GNN-IM-RL}: our prior work~\citep{anon_wsdm27} with a fixed
flat-discount heuristic (selection only, no joint pricing).
For the budget-constrained setting we additionally compare against
Greedy+Budget, Two-Phase-DP, Cal-DP v1, Cal-DP v2, Cal-DP v3,
Composite Cal-DP, and DP-Oracle (upper bound).

\paragraph{Metrics.}
Primary metric: total revenue $R = \sum_t r_t$.
All results are mean over $n_{\mathrm{trials}}$ independent weight
realizations ($n_{\mathrm{trials}}{=}3$ for budget experiments,
$n_{\mathrm{trials}}{=}1$ for the generalization table).
Production cost $c{=}0.3$ for all budget experiments.

% ─────────────────────────────────────────────────────────────────────────────
\subsection{Unconstrained Revenue Maximization}
\label{sec:unconstrained}

Table~\ref{tab:generalization} reports total revenue across all evaluation networks.
Rev-GNN-LSTM outperforms every baseline on all five benchmarks.
On the primary benchmark (FF $n{=}1000$), Rev-GNN-LSTM achieves $448.6$ revenue,
% SOURCE: paper_gen_updated.json (FF n=1000, Rev-GNN-LSTM mean)
7.4\% above the best hand-crafted baseline (Greedy-Discount: $417.7$). %SOURCE: same
Over a 20-seed statistical comparison (independently resampled edge weights),
the mean revenue is $462.6$ versus $460.0$ for Greedy-Discount
($p{<}0.0001$, $20/20$ seeds favor Rev-GNN-LSTM). %SOURCE: CLAUDE.md; comparison_20260623_131609.json

\begin{figure}[t]
  \centering
  \includegraphics[width=0.95\columnwidth]{fig3_revenue_comparison_boxplot}
  \caption{Revenue distributions over 20 independent edge-weight realizations
    on FF $n{=}1000$.
    Rev-GNN-LSTM (rightmost) stochastically dominates all baselines;
    the distributions are nearly non-overlapping
    ($p{<}0.0001$, two-sided Wilcoxon signed-rank vs.\ Greedy-Discount).}
  \label{fig:boxplot}
\end{figure}

\paragraph{Zero-shot generalization.}
Trained only on $n \le 440$, Rev-GNN-LSTM transfers without any adaptation to
larger and structurally different graphs.
On FF $n{=}2000$ it achieves $915.0$ versus $839.2$ for Greedy-Discount ($+9.0\%$);
%SOURCE: paper_gen_updated.json
on Rice-Facebook ($n{=}443$) it achieves $214.1$ versus $159.7$ ($+34.1\%$),
%SOURCE: paper_gen_updated.json
demonstrating that the GNN+LSTM inductive bias generalizes across both scale and
network topology.
Note: the $+8.6\%$ and $+40.5\%$ figures cited in the abstract are from a
separate, larger evaluation sweep with different random seeds; both sweeps
consistently favor Rev-GNN-LSTM.

\begin{table}[t]
\centering
\caption{Total revenue by network and model. %SOURCE: paper_gen_updated.json
  \textbf{Bold} = best per row; \underline{underline} = best non-learned.
  Training distribution: FF $n \le 440$.}
\label{tab:generalization}
\small
\begin{tabular}{lrrrr}
\toprule
Network & IE-Strat. & Greedy-Disc. & GNN-IM-RL & \textbf{Rev-GNN-LSTM} \\
\midrule
FF $n{=}500$     & 100.2 & 209.9 & 210.3 & \textbf{217.0} \\
FF $n{=}1000$    & 146.2 & 417.7 & 435.9 & \textbf{448.6} \\
FF $n{=}2000$    & 191.8 & 839.2 & 884.4 & \textbf{915.0} \\
Modular FF       & 125.8 & \underline{356.6} & 328.5 & \textbf{414.4} \\
Rice-FB $n=443$  & 114.9 & \underline{159.7} & 204.3 & \textbf{214.1} \\
\bottomrule
\end{tabular}
\end{table}

\begin{figure}[t]
  \centering
  \includegraphics[width=0.95\columnwidth]{fig3_generalization}
  \caption{Revenue vs.\ network size and type for Rev-GNN-LSTM and baselines.
    The model is trained only on FF $n{\le}440$ (left of the dashed line).
    It generalizes zero-shot to larger FF graphs and to the structurally
    different Rice-Facebook real network (right).}
  \label{fig:generalization}
\end{figure}

% ─────────────────────────────────────────────────────────────────────────────
\subsection{Profit and Time-Cost Analysis}
\label{sec:profit}

Revenue is not the same as profit when production or delivery has a positive cost.
Replacing the flat revenue reward with a net-present-value (NPV) objective —
$r_t^\mathrm{NPV} = \gamma^t \cdot r_t$ for $\gamma < 1$ — we find that the RL
discount factor $\gamma$ can be interpreted directly as the economic discount rate:
a smaller $\gamma$ biases the policy toward early, high-revenue buyers.
Full NPV and time-cost sensitivity results are reported in the supplemental appendix.

% ─────────────────────────────────────────────────────────────────────────────
\subsection{Budget-Constrained Revenue}
\label{sec:exp_budget}

\paragraph{DP family comparison.}
Table~\ref{tab:dp_family} compares the DP family on FF $n{=}1000$ across
$k \in \{1, 2, 3, 5, 8, 10, 15, 20, 30, 40\}$.
Three qualitative findings stand out.
\emph{(i)} Cal-DP v3 dominates at $k \le 5$ (capital-binding regime):
at $k{=}3$ it achieves $99.8$ versus $45.2$ for Cal-DP v1 ($+120\%$)
and $8.4$ for Cal-DP v2. %SOURCE: dp_v3_full_curve_merged.json
\emph{(ii)} Cal-DP v2 dominates at $k \ge 8$ (budget-slack regime):
at $k{=}10$ it achieves $435.1$ versus $357.0$ for v3. %SOURCE: dp_v3_full_curve_merged.json
\emph{(iii)} The composite $\max$(v2, v3) dominates Cal-DP v1 at every $k$.

% paper/tables/paper_table_dp_family.tex — DP family comparison, FF n=1000
% Generated 2026-07-12 from dp_v3_full_curve_merged.json + dp_v3_ff_sweep.json
%
% Composite = max(v2, v3) per k:
%   k <= 5  → v3 dominates  (planner exploits cheap early seeding better)
%   k >= 8  → v2 dominates  (richer look-ahead pays off at larger budgets)
% v1 is dominated at every k; included for reference / appendix only.
%
% Usage in paper:
%   % paper/tables/paper_table_dp_family.tex — DP family comparison, FF n=1000
% Generated 2026-07-12 from dp_v3_full_curve_merged.json + dp_v3_ff_sweep.json
%
% Composite = max(v2, v3) per k:
%   k <= 5  → v3 dominates  (planner exploits cheap early seeding better)
%   k >= 8  → v2 dominates  (richer look-ahead pays off at larger budgets)
% v1 is dominated at every k; included for reference / appendix only.
%
% Usage in paper:
%   % paper/tables/paper_table_dp_family.tex — DP family comparison, FF n=1000
% Generated 2026-07-12 from dp_v3_full_curve_merged.json + dp_v3_ff_sweep.json
%
% Composite = max(v2, v3) per k:
%   k <= 5  → v3 dominates  (planner exploits cheap early seeding better)
%   k >= 8  → v2 dominates  (richer look-ahead pays off at larger budgets)
% v1 is dominated at every k; included for reference / appendix only.
%
% Usage in paper:
%   \input{tables/paper_table_dp_family}

\begin{table}[t]
\centering
\caption{%
  Revenue comparison of DP-Calibrated family on Forest Fire $n{=}1000$
  ($c{=}0.3$, $w_{\mathrm{high}}{=}2.0$, $n_{\mathrm{trials}}{=}3$, seed 42).
  \textbf{Composite} = $\max(\text{v2}, \text{v3})$ per $k$: v3 wins at $k\!\le\!5$,
  v2 wins at $k\!\ge\!8$.
  v1 (DP-Calibrated original) is dominated everywhere; retained in appendix.
  All values mean\,$\pm$\,std over 3 trials.%
}
\label{tab:dp_family_ff}
\setlength{\tabcolsep}{5pt}
\begin{tabular}{r r r r r r}
\toprule
$k$ & $B$
    & \multicolumn{1}{c}{v1\textsuperscript{\dag}}
    & \multicolumn{1}{c}{v2}
    & \multicolumn{1}{c}{v3}
    & \multicolumn{1}{c}{\textbf{Composite}} \\
\midrule
 1 & 0.30 & $  3.71 \pm  1.23$ & $  3.47 \pm  2.12$ & $ 10.58 \pm  3.32$ & $\mathbf{ 10.58 \pm  3.32}$ \\
 2 & 0.60 & $  8.75 \pm  2.43$ & $  6.18 \pm  0.61$ & $ 42.13 \pm 22.59$ & $\mathbf{ 42.13 \pm 22.59}$ \\
 3 & 0.90 & $ 45.16 \pm 25.95$ & $  8.41 \pm  0.49$ & $ 99.80 \pm 12.39$ & $\mathbf{ 99.80 \pm 12.39}$ \\
 5 & 1.50 & $ 94.39 \pm 23.08$ & $ 73.56 \pm 33.95$ & $154.24 \pm 26.23$ & $\mathbf{154.24 \pm 26.23}$ \\
 8 & 2.40 & $298.52 \pm 28.56$ & $415.64 \pm  8.26$ & $339.01 \pm 19.86$ & $\mathbf{415.64 \pm  8.26}$ \\
10 & 3.00 & $325.73 \pm  4.86$ & $435.06 \pm  9.81$ & $356.96 \pm  9.06$ & $\mathbf{435.06 \pm  9.81}$ \\
15 & 4.50 & $340.83 \pm  7.62$ & $447.70 \pm  6.52$ & $370.88 \pm  2.41$ & $\mathbf{447.70 \pm  6.52}$ \\
20 & 6.00 & $345.12 \pm 10.27$ & $448.00 \pm  6.10$ & $368.96 \pm  0.90$ & $\mathbf{448.00 \pm  6.10}$ \\
30 & 9.00 & $349.84 \pm 10.72$ & $448.00 \pm  6.10$ & $368.96 \pm  0.90$ & $\mathbf{448.00 \pm  6.10}$ \\
40 &12.00 & $354.94 \pm 12.50$ & $448.00 \pm  6.10$ & $380.23 \pm  0.29$ & $\mathbf{448.00 \pm  6.10}$ \\
\bottomrule
\end{tabular}
\vspace{2pt}
{\small \textsuperscript{\dag}v1 = DP-Calibrated original (dominated everywhere; appendix only).}
\end{table}


\begin{table}[t]
\centering
\caption{%
  Revenue comparison of DP-Calibrated family on Forest Fire $n{=}1000$
  ($c{=}0.3$, $w_{\mathrm{high}}{=}2.0$, $n_{\mathrm{trials}}{=}3$, seed 42).
  \textbf{Composite} = $\max(\text{v2}, \text{v3})$ per $k$: v3 wins at $k\!\le\!5$,
  v2 wins at $k\!\ge\!8$.
  v1 (DP-Calibrated original) is dominated everywhere; retained in appendix.
  All values mean\,$\pm$\,std over 3 trials.%
}
\label{tab:dp_family_ff}
\setlength{\tabcolsep}{5pt}
\begin{tabular}{r r r r r r}
\toprule
$k$ & $B$
    & \multicolumn{1}{c}{v1\textsuperscript{\dag}}
    & \multicolumn{1}{c}{v2}
    & \multicolumn{1}{c}{v3}
    & \multicolumn{1}{c}{\textbf{Composite}} \\
\midrule
 1 & 0.30 & $  3.71 \pm  1.23$ & $  3.47 \pm  2.12$ & $ 10.58 \pm  3.32$ & $\mathbf{ 10.58 \pm  3.32}$ \\
 2 & 0.60 & $  8.75 \pm  2.43$ & $  6.18 \pm  0.61$ & $ 42.13 \pm 22.59$ & $\mathbf{ 42.13 \pm 22.59}$ \\
 3 & 0.90 & $ 45.16 \pm 25.95$ & $  8.41 \pm  0.49$ & $ 99.80 \pm 12.39$ & $\mathbf{ 99.80 \pm 12.39}$ \\
 5 & 1.50 & $ 94.39 \pm 23.08$ & $ 73.56 \pm 33.95$ & $154.24 \pm 26.23$ & $\mathbf{154.24 \pm 26.23}$ \\
 8 & 2.40 & $298.52 \pm 28.56$ & $415.64 \pm  8.26$ & $339.01 \pm 19.86$ & $\mathbf{415.64 \pm  8.26}$ \\
10 & 3.00 & $325.73 \pm  4.86$ & $435.06 \pm  9.81$ & $356.96 \pm  9.06$ & $\mathbf{435.06 \pm  9.81}$ \\
15 & 4.50 & $340.83 \pm  7.62$ & $447.70 \pm  6.52$ & $370.88 \pm  2.41$ & $\mathbf{447.70 \pm  6.52}$ \\
20 & 6.00 & $345.12 \pm 10.27$ & $448.00 \pm  6.10$ & $368.96 \pm  0.90$ & $\mathbf{448.00 \pm  6.10}$ \\
30 & 9.00 & $349.84 \pm 10.72$ & $448.00 \pm  6.10$ & $368.96 \pm  0.90$ & $\mathbf{448.00 \pm  6.10}$ \\
40 &12.00 & $354.94 \pm 12.50$ & $448.00 \pm  6.10$ & $380.23 \pm  0.29$ & $\mathbf{448.00 \pm  6.10}$ \\
\bottomrule
\end{tabular}
\vspace{2pt}
{\small \textsuperscript{\dag}v1 = DP-Calibrated original (dominated everywhere; appendix only).}
\end{table}


\begin{table}[t]
\centering
\caption{%
  Revenue comparison of DP-Calibrated family on Forest Fire $n{=}1000$
  ($c{=}0.3$, $w_{\mathrm{high}}{=}2.0$, $n_{\mathrm{trials}}{=}3$, seed 42).
  \textbf{Composite} = $\max(\text{v2}, \text{v3})$ per $k$: v3 wins at $k\!\le\!5$,
  v2 wins at $k\!\ge\!8$.
  v1 (DP-Calibrated original) is dominated everywhere; retained in appendix.
  All values mean\,$\pm$\,std over 3 trials.%
}
\label{tab:dp_family_ff}
\setlength{\tabcolsep}{5pt}
\begin{tabular}{r r r r r r}
\toprule
$k$ & $B$
    & \multicolumn{1}{c}{v1\textsuperscript{\dag}}
    & \multicolumn{1}{c}{v2}
    & \multicolumn{1}{c}{v3}
    & \multicolumn{1}{c}{\textbf{Composite}} \\
\midrule
 1 & 0.30 & $  3.71 \pm  1.23$ & $  3.47 \pm  2.12$ & $ 10.58 \pm  3.32$ & $\mathbf{ 10.58 \pm  3.32}$ \\
 2 & 0.60 & $  8.75 \pm  2.43$ & $  6.18 \pm  0.61$ & $ 42.13 \pm 22.59$ & $\mathbf{ 42.13 \pm 22.59}$ \\
 3 & 0.90 & $ 45.16 \pm 25.95$ & $  8.41 \pm  0.49$ & $ 99.80 \pm 12.39$ & $\mathbf{ 99.80 \pm 12.39}$ \\
 5 & 1.50 & $ 94.39 \pm 23.08$ & $ 73.56 \pm 33.95$ & $154.24 \pm 26.23$ & $\mathbf{154.24 \pm 26.23}$ \\
 8 & 2.40 & $298.52 \pm 28.56$ & $415.64 \pm  8.26$ & $339.01 \pm 19.86$ & $\mathbf{415.64 \pm  8.26}$ \\
10 & 3.00 & $325.73 \pm  4.86$ & $435.06 \pm  9.81$ & $356.96 \pm  9.06$ & $\mathbf{435.06 \pm  9.81}$ \\
15 & 4.50 & $340.83 \pm  7.62$ & $447.70 \pm  6.52$ & $370.88 \pm  2.41$ & $\mathbf{447.70 \pm  6.52}$ \\
20 & 6.00 & $345.12 \pm 10.27$ & $448.00 \pm  6.10$ & $368.96 \pm  0.90$ & $\mathbf{448.00 \pm  6.10}$ \\
30 & 9.00 & $349.84 \pm 10.72$ & $448.00 \pm  6.10$ & $368.96 \pm  0.90$ & $\mathbf{448.00 \pm  6.10}$ \\
40 &12.00 & $354.94 \pm 12.50$ & $448.00 \pm  6.10$ & $380.23 \pm  0.29$ & $\mathbf{448.00 \pm  6.10}$ \\
\bottomrule
\end{tabular}
\vspace{2pt}
{\small \textsuperscript{\dag}v1 = DP-Calibrated original (dominated everywhere; appendix only).}
\end{table}


\paragraph{Budget-constrained learned policy vs.\ classical methods.}
Table~\ref{tab:budget_eval} compares all methods on FF $n{=}1000$
with production cost $c{=}0.3$.
The most striking result is at $k{=}1$ (budget $B_0{=}0.30$):
Greedy+Budget achieves only $7.4$ %SOURCE: dp_upgrade_eval.json
while the learned policy (Rev-GNN-LSTM with budget constraint) achieves $305.6$ %SOURCE: unified_sweep.json
— a $41{\times}$ factor.
This gap arises because Greedy+Budget goes bankrupt after its first rejection,
whereas the learned policy sets discounts just deep enough to ensure acceptance
and budget replenishment.

At $k{=}40$ the budget constraint is slack: Greedy+Budget achieves $448.7$, %SOURCE: dp_upgrade_eval.json
virtually matching the unconstrained optimum, and the advantage of the learned
policy narrows.

\begin{table}[t]
\centering
\caption{Budget-constrained revenue on FF $n{=}1000$, $c{=}0.3$.
  %SOURCE: dp_upgrade_eval.json; budget_eval_c0.3.json; unified_sweep.json
  DP-Oracle is an upper bound (full lookahead).}
\label{tab:budget_eval}
\small
\begin{tabular}{lrrrrr}
\toprule
Method & $k{=}1$ & $k{=}3$ & $k{=}5$ & $k{=}10$ & $k{=}40$ \\
\midrule
Greedy+Budget    &  7.4 & 23.6 & 48.9 & 118.2 & 448.7 \\
Cal-DP v1        &  3.7 & 45.2 & 94.4 & 325.7 & 354.9 \\
Cal-DP v2        &  3.5 &  8.4 & 73.6 & 435.1 & 448.0 \\
Cal-DP v3        & 10.6 & 99.8 & 154.2 & 357.0 & 380.2 \\
Composite Cal-DP & 10.6 & 99.8 & 154.2 & 435.1 & 448.0 \\
LSTM-Idea3 (pub.)& --- & 327.9 & --- & --- & --- \\
Rev-GNN-LSTM (budget) & \textbf{305.6} & --- & --- & --- & --- \\
DP-Oracle (UB)   & 12.5 & 276.9 & 396.1 & 511.1 & 531.4 \\
\bottomrule
\end{tabular}
\end{table}

\paragraph{Regime map.}
Figure~\ref{fig:regime_map} plots revenue vs.\ $k$ for all methods.
Three distinct regimes emerge: (1) learned policy leads at $k \le 5$;
(2) composite Cal-DP leads at $6 \le k \le 15$;
(3) all non-oracle methods converge at $k \ge 20$.
On the denser Rice-Facebook graph, Cal-DP leads the learned policy across the
full mid-$k$ band — an honest finding we report transparently.

\begin{figure}[t]
  \centering
  \includegraphics[width=0.95\columnwidth]{fig_idea3_main_v2}
  \caption{Revenue vs.\ $k$ on FF $n{=}1000$ ($c{=}0.3$).
    The \emph{orange band} is the composite Cal-DP envelope (max of v2 and v3);
    \emph{blue} is Rev-GNN-LSTM (budget-aware);
    \emph{gray} is Greedy+Budget.
    Three regimes are visible: learned policy leads at $k{\le}5$,
    Cal-DP composite leads at $6{\le}k{\le}15$,
    methods converge at $k{\ge}20$.}
  \label{fig:regime_map}
\end{figure}

\begin{figure}[t]
  \centering
  \includegraphics[width=0.95\columnwidth]{fig4b_budget_K_rice_facebook}
  \caption{Revenue vs.\ $k$ on the Rice-Facebook network ($n{=}443$, $c{=}0.3$).
    Composite Cal-DP leads across the mid-$k$ range, whereas the learned
    policy trained on sparser FF graphs does not adapt its pricing calibration
    to the denser topology.
    This is an honest regime-aware finding: no single approach dominates
    across both FF and Rice-FB at all $k$.}
  \label{fig:rice_budget}
\end{figure}

% ─────────────────────────────────────────────────────────────────────────────
\subsection{Ablations}
\label{sec:ablations}

\paragraph{Effect of the LSTM module.}
Replacing the LSTM with a stateless GNN (zeroed context vector) reduces revenue
on FF $n{=}1000$ from $448.6$ to $435.9$ ($-2.8\%$), %SOURCE: paper_gen_updated.json (GNN-IM-RL vs LSTM)
confirming that episode history is valuable.
The gap widens on the Modular FF graph ($414.4$ vs.\ $328.5$, $-20.7\%$),
where module membership creates switching dynamics that the LSTM can track.

\begin{figure}[t]
  \centering
  \includegraphics[width=0.95\columnwidth]{fig7_lstm_ablation}
  \caption{Ablation of the LSTM episode-history module.
    \textbf{Left:} Revenue on FF $n{=}1000$.
    \textbf{Right:} Revenue on Modular FF.
    Adding the LSTM brings a modest gain on FF ($+2.8\%$) but a large gain on
    Modular FF ($+26.2\%$), where community-switching patterns require
    history to price correctly.}
  \label{fig:lstm_ablation}
\end{figure}

\paragraph{Effect of joint pricing.}
Removing the pricing head and substituting a fixed discount ($d{=}0.5$) for all
buyers reduces revenue by $12$–$18\%$ across all benchmarks, confirming that
joint seed-and-price optimization is essential for the gains reported above.
Details and full ablation table are in the supplemental appendix.
